% --- Archon physics formalization source begin ---
% archon:physics
% archon:covers PhyXMiniProblems/problem_phyx_mini_0392.lean
% archon:source-report reports/phyx_mini/problem_phyx_mini_0392.source.json

\chapter{Physics problem phyx\_mini\_0392}
\label{ch:PhyXMiniProblems_problem_phyx_mini_0392}

\paragraph{Problem source.}
\#\# Physical scenario

A dam retains a lake 6 m deep. To construct a gate in the dam, we need to know the net horizontal force on a 5-m-wide, 6-m-tall port section that then replaces a 5-m section of the dam.

\#\# Question

Find the net horizontal force from the water on one side and air on the other side of the port.

\#\# Answer choices

- A: A: 490kN

- B: B: 880kN

- C: C: 154kN

- D: D: 10.2kN

\#\# Figure caption (auxiliary; use the image as primary evidence)

The image consists of two diagrams showing a side view and a top view of a lake.

**Side View:**
- A body of water labeled "Lake" is depicted with a horizontal water surface.
- A solid structure is shown with vertical and sloped lines next to the lake.
- The height of the lake level above a reference point is labeled as "6 m".
- The structure appears to have a supporting or retaining function.

**Top View:**
- The lake is shown as a horizontal expanse, labeled "Lake".
- Two thick, horizontal lines represent land or solid boundaries on either side of the lake.
- A narrow, rectangular bridge or structure connecting the two land boundaries over the lake is shown.
- The width of the bridge or structure is labeled as "5 m".

Both diagrams provide contextual information about the relative positioning and dimensions of the structures in relation to the lake.

\#\# Dataset metadata

Laws of Thermodynamics; Physical Model Grounding Reasoning, Multi-Formula Reasoning

\paragraph{Recorded answer/context.}
B: B: 880kN

\paragraph{Figure/image path.}
/root/proposal\_for\_physic/science-mango/phyx\_mini\_run/phyx\_data/test\_image/392.png

\paragraph{Formalization target.}
create a compiling Lean file with sorry bodies at `PhyXMiniProblems/problem\_phyx\_mini\_0392.lean`. The Lean declarations must preserve the physical quantities, dimensions or dimensional roles, figure labels, governing-law hypotheses, and final relation expressed by this problem.
Use Mathlib/Physlib names found through LeanExplore where available. If a physics API is missing, introduce faithful local abstractions rather than scalar placeholder aliases.

\begin{theorem}[Physics formalization target]
\label{thm:physics:phyx_mini_0392:target}
\lean{PhyXMiniProblems.ProblemPhyXMini0392.problem_phyx_mini_0392}
\uses{def:physics:phyx-mini-0392:phyxminiproblems-problemphyxmini0392-forceinkilonewtons, def:physics:phyx-mini-0392:phyxminiproblems-problemphyxmini0392-damportsetup, def:physics:phyx-mini-0392:phyxminiproblems-problemphyxmini0392-matchesstatedproblemdata, def:physics:phyx-mini-0392:phyxminiproblems-problemphyxmini0392-matchesprimarydamportfigure, def:physics:phyx-mini-0392:phyxminiproblems-problemphyxmini0392-hasstandardwatergravitycalibration, def:physics:phyx-mini-0392:phyxminiproblems-problemphyxmini0392-satisfieshydrostaticpressurelaw, def:physics:phyx-mini-0392:phyxminiproblems-problemphyxmini0392-satisfiespressureresultantlaw, def:physics:phyx-mini-0392:phyxminiproblems-problemphyxmini0392-isuniqueclosestanswer}
The assigned autoformalize agent should translate this physics problem into Lean declarations in the covered file, with theorem and lemma proof bodies written as `by sorry`.
\end{theorem}
\begin{proof}
This is an autoformalization task, not a proof task. Produce statements that can later be proved by the physics prover without weakening the physical model.
\end{proof}

% --- Archon declaration topology begin ---
\section{Declaration topology}
\begin{definition}[mass Density Dimension]
  \label{def:physics:phyx-mini-0392:phyxminiproblems-problemphyxmini0392-massdensitydimension}
  \lean{PhyXMiniProblems.ProblemPhyXMini0392.massDensityDimension}
  The physical dimension of mass density, M L⁻³
\end{definition}
\begin{proof}
  This is the declared model definition of mass Density Dimension.
\end{proof}
\begin{definition}[acceleration Dimension]
  \label{def:physics:phyx-mini-0392:phyxminiproblems-problemphyxmini0392-accelerationdimension}
  \lean{PhyXMiniProblems.ProblemPhyXMini0392.accelerationDimension}
  The physical dimension of acceleration, L T⁻²
\end{definition}
\begin{proof}
  This is the declared model definition of acceleration Dimension.
\end{proof}
\begin{definition}[force Dimension]
  \label{def:physics:phyx-mini-0392:phyxminiproblems-problemphyxmini0392-forcedimension}
  \lean{PhyXMiniProblems.ProblemPhyXMini0392.forceDimension}
  The physical dimension of force, M L T⁻²
\end{definition}
\begin{proof}
  This is the declared model definition of force Dimension.
\end{proof}
\begin{definition}[Length Quantity]
  \label{def:physics:phyx-mini-0392:phyxminiproblems-problemphyxmini0392-lengthquantity}
  \lean{PhyXMiniProblems.ProblemPhyXMini0392.LengthQuantity}
  A nonnegative physical length independent of the chosen readout unit
\end{definition}
\begin{proof}
  This is the declared model definition of Length Quantity.
\end{proof}
\begin{definition}[Mass Density Quantity]
  \label{def:physics:phyx-mini-0392:phyxminiproblems-problemphyxmini0392-massdensityquantity}
  \lean{PhyXMiniProblems.ProblemPhyXMini0392.MassDensityQuantity}
  \uses{def:physics:phyx-mini-0392:phyxminiproblems-problemphyxmini0392-massdensitydimension}
  A nonnegative physical mass density
\end{definition}
\begin{proof}
  This is the declared model definition of Mass Density Quantity.
\end{proof}
\begin{definition}[Acceleration Quantity]
  \label{def:physics:phyx-mini-0392:phyxminiproblems-problemphyxmini0392-accelerationquantity}
  \lean{PhyXMiniProblems.ProblemPhyXMini0392.AccelerationQuantity}
  \uses{def:physics:phyx-mini-0392:phyxminiproblems-problemphyxmini0392-accelerationdimension}
  A nonnegative physical acceleration magnitude
\end{definition}
\begin{proof}
  This is the declared model definition of Acceleration Quantity.
\end{proof}
\begin{definition}[Force Quantity]
  \label{def:physics:phyx-mini-0392:phyxminiproblems-problemphyxmini0392-forcequantity}
  \lean{PhyXMiniProblems.ProblemPhyXMini0392.ForceQuantity}
  \uses{def:physics:phyx-mini-0392:phyxminiproblems-problemphyxmini0392-forcedimension}
  A nonnegative physical force magnitude
\end{definition}
\begin{proof}
  This is the declared model definition of Force Quantity.
\end{proof}
\begin{definition}[Pressure Quantity]
  \label{def:physics:phyx-mini-0392:phyxminiproblems-problemphyxmini0392-pressurequantity}
  \lean{PhyXMiniProblems.ProblemPhyXMini0392.PressureQuantity}
  A physical absolute pressure, using Physlib's pressure dimension
\end{definition}
\begin{proof}
  This is the declared model definition of Pressure Quantity.
\end{proof}
\begin{definition}[length In Meters]
  \label{def:physics:phyx-mini-0392:phyxminiproblems-problemphyxmini0392-lengthinmeters}
  \lean{PhyXMiniProblems.ProblemPhyXMini0392.lengthInMeters}
  \uses{def:physics:phyx-mini-0392:phyxminiproblems-problemphyxmini0392-lengthquantity}
  Read a physical length in SI metres
\end{definition}
\begin{proof}
  This is the declared model definition of length In Meters.
\end{proof}
\begin{definition}[mass Density In Kilograms Per Cubic Meter]
  \label{def:physics:phyx-mini-0392:phyxminiproblems-problemphyxmini0392-massdensityinkilogramspercubicmeter}
  \lean{PhyXMiniProblems.ProblemPhyXMini0392.massDensityInKilogramsPerCubicMeter}
  \uses{def:physics:phyx-mini-0392:phyxminiproblems-problemphyxmini0392-massdensityquantity}
  Read a physical mass density in SI kilograms per cubic metre
\end{definition}
\begin{proof}
  This is the declared model definition of mass Density In Kilograms Per Cubic Meter.
\end{proof}
\begin{definition}[acceleration In Meters Per Second Squared]
  \label{def:physics:phyx-mini-0392:phyxminiproblems-problemphyxmini0392-accelerationinmeterspersecondsquared}
  \lean{PhyXMiniProblems.ProblemPhyXMini0392.accelerationInMetersPerSecondSquared}
  \uses{def:physics:phyx-mini-0392:phyxminiproblems-problemphyxmini0392-accelerationquantity}
  Read a physical acceleration in SI metres per second squared
\end{definition}
\begin{proof}
  This is the declared model definition of acceleration In Meters Per Second Squared.
\end{proof}
\begin{definition}[pressure In Pascals]
  \label{def:physics:phyx-mini-0392:phyxminiproblems-problemphyxmini0392-pressureinpascals}
  \lean{PhyXMiniProblems.ProblemPhyXMini0392.pressureInPascals}
  \uses{def:physics:phyx-mini-0392:phyxminiproblems-problemphyxmini0392-pressurequantity}
  Read a physical absolute pressure in SI pascals
\end{definition}
\begin{proof}
  This is the declared model definition of pressure In Pascals.
\end{proof}
\begin{definition}[force In Newtons]
  \label{def:physics:phyx-mini-0392:phyxminiproblems-problemphyxmini0392-forceinnewtons}
  \lean{PhyXMiniProblems.ProblemPhyXMini0392.forceInNewtons}
  \uses{def:physics:phyx-mini-0392:phyxminiproblems-problemphyxmini0392-forcequantity}
  Read a physical force magnitude in SI newtons
\end{definition}
\begin{proof}
  This is the declared model definition of force In Newtons.
\end{proof}
\begin{definition}[force In Kilonewtons]
  \label{def:physics:phyx-mini-0392:phyxminiproblems-problemphyxmini0392-forceinkilonewtons}
  \lean{PhyXMiniProblems.ProblemPhyXMini0392.forceInKilonewtons}
  \uses{def:physics:phyx-mini-0392:phyxminiproblems-problemphyxmini0392-forcequantity, def:physics:phyx-mini-0392:phyxminiproblems-problemphyxmini0392-forceinnewtons}
  Read a physical force magnitude in kilonewtons
\end{definition}
\begin{proof}
  This is the declared model definition of force In Kilonewtons.
\end{proof}
\begin{definition}[Port Side]
  \label{def:physics:phyx-mini-0392:phyxminiproblems-problemphyxmini0392-portside}
  \lean{PhyXMiniProblems.ProblemPhyXMini0392.PortSide}
  The two faces of the port distinguished by the scenario
\end{definition}
\begin{proof}
  This is the declared model definition of Port Side.
\end{proof}
\begin{definition}[Fluid]
  \label{def:physics:phyx-mini-0392:phyxminiproblems-problemphyxmini0392-fluid}
  \lean{PhyXMiniProblems.ProblemPhyXMini0392.Fluid}
  Fluids acting on the two faces of the port
\end{definition}
\begin{proof}
  This is the declared model definition of Fluid.
\end{proof}
\begin{definition}[Port Orientation]
  \label{def:physics:phyx-mini-0392:phyxminiproblems-problemphyxmini0392-portorientation}
  \lean{PhyXMiniProblems.ProblemPhyXMini0392.PortOrientation}
  Orientation of the rectangular port relative to the lake surface
\end{definition}
\begin{proof}
  This is the declared model definition of Port Orientation.
\end{proof}
\begin{definition}[Figure Panel]
  \label{def:physics:phyx-mini-0392:phyxminiproblems-problemphyxmini0392-figurepanel}
  \lean{PhyXMiniProblems.ProblemPhyXMini0392.FigurePanel}
  The two explicitly named views in the supplied raster
\end{definition}
\begin{proof}
  This is the declared model definition of Figure Panel.
\end{proof}
\begin{definition}[Figure Text Label]
  \label{def:physics:phyx-mini-0392:phyxminiproblems-problemphyxmini0392-figuretextlabel}
  \lean{PhyXMiniProblems.ProblemPhyXMini0392.FigureTextLabel}
  Literal text labels visible in the supplied raster
\end{definition}
\begin{proof}
  This is the declared model definition of Figure Text Label.
\end{proof}
\begin{definition}[Figure Length Mark]
  \label{def:physics:phyx-mini-0392:phyxminiproblems-problemphyxmini0392-figurelengthmark}
  \lean{PhyXMiniProblems.ProblemPhyXMini0392.FigureLengthMark}
  The two dimensionful arrows shown in the supplied raster
\end{definition}
\begin{proof}
  This is the declared model definition of Figure Length Mark.
\end{proof}
\begin{definition}[Dam Port Figure]
  \label{def:physics:phyx-mini-0392:phyxminiproblems-problemphyxmini0392-damportfigure}
  \lean{PhyXMiniProblems.ProblemPhyXMini0392.DamPortFigure}
  \uses{def:physics:phyx-mini-0392:phyxminiproblems-problemphyxmini0392-lengthquantity, def:physics:phyx-mini-0392:phyxminiproblems-problemphyxmini0392-figurepanel, def:physics:phyx-mini-0392:phyxminiproblems-problemphyxmini0392-figuretextlabel, def:physics:phyx-mini-0392:phyxminiproblems-problemphyxmini0392-figurelengthmark}
  Qualitative and dimensionful evidence transcribed from the primary image. Its 6 m arrow measures the lake depth in the side view, while its 5 m arrow measures the port opening in the top view. No force value is stored here
\end{definition}
\begin{proof}
  This is the declared model definition of Dam Port Figure.
\end{proof}
\begin{definition}[Dam Port Setup]
  \label{def:physics:phyx-mini-0392:phyxminiproblems-problemphyxmini0392-damportsetup}
  \lean{PhyXMiniProblems.ProblemPhyXMini0392.DamPortSetup}
  \uses{def:physics:phyx-mini-0392:phyxminiproblems-problemphyxmini0392-lengthquantity, def:physics:phyx-mini-0392:phyxminiproblems-problemphyxmini0392-massdensityquantity, def:physics:phyx-mini-0392:phyxminiproblems-problemphyxmini0392-accelerationquantity, def:physics:phyx-mini-0392:phyxminiproblems-problemphyxmini0392-forcequantity, def:physics:phyx-mini-0392:phyxminiproblems-problemphyxmini0392-pressurequantity, def:physics:phyx-mini-0392:phyxminiproblems-problemphyxmini0392-portside, def:physics:phyx-mini-0392:phyxminiproblems-problemphyxmini0392-fluid, def:physics:phyx-mini-0392:phyxminiproblems-problemphyxmini0392-portorientation, def:physics:phyx-mini-0392:phyxminiproblems-problemphyxmini0392-damportfigure}
  Independent quantities in the dam-port experiment. The pressure field takes a scalar argument because that argument is explicitly a metre readout of the depth measured downward from the lake surface; each field value is still a dimensionful physical pressure. The requested force is an independent observable constrained only by the governing law below
\end{definition}
\begin{proof}
  This is the declared model definition of Dam Port Setup.
\end{proof}
\begin{definition}[Matches Stated Problem Data]
  \label{def:physics:phyx-mini-0392:phyxminiproblems-problemphyxmini0392-matchesstatedproblemdata}
  \lean{PhyXMiniProblems.ProblemPhyXMini0392.MatchesStatedProblemData}
  \uses{def:physics:phyx-mini-0392:phyxminiproblems-problemphyxmini0392-lengthinmeters, def:physics:phyx-mini-0392:phyxminiproblems-problemphyxmini0392-damportsetup}
  Numerical and qualitative data stated in the prose problem
\end{definition}
\begin{proof}
  This is the declared model definition of Matches Stated Problem Data.
\end{proof}
\begin{definition}[Matches Primary Dam Port Figure]
  \label{def:physics:phyx-mini-0392:phyxminiproblems-problemphyxmini0392-matchesprimarydamportfigure}
  \lean{PhyXMiniProblems.ProblemPhyXMini0392.MatchesPrimaryDamPortFigure}
  \uses{def:physics:phyx-mini-0392:phyxminiproblems-problemphyxmini0392-damportsetup}
  Primary-image evidence: both named views and all printed labels are visible; the side-view arrow is the lake depth and the top-view arrow is the port width. These observations contain no pressure or force answer
\end{definition}
\begin{proof}
  This is the declared model definition of Matches Primary Dam Port Figure.
\end{proof}
\begin{definition}[Has Standard Water Gravity Calibration]
  \label{def:physics:phyx-mini-0392:phyxminiproblems-problemphyxmini0392-hasstandardwatergravitycalibration}
  \lean{PhyXMiniProblems.ProblemPhyXMini0392.HasStandardWaterGravityCalibration}
  \uses{def:physics:phyx-mini-0392:phyxminiproblems-problemphyxmini0392-massdensityinkilogramspercubicmeter, def:physics:phyx-mini-0392:phyxminiproblems-problemphyxmini0392-accelerationinmeterspersecondsquared, def:physics:phyx-mini-0392:phyxminiproblems-problemphyxmini0392-pressureinpascals, def:physics:phyx-mini-0392:phyxminiproblems-problemphyxmini0392-damportsetup}
  The standard textbook calibration implicit in the recorded answer: water is read as 1000 kg/m³ and gravitational acceleration as 9.8 m/s². Ambient pressure is positive but need not be assigned a numerical value because it cancels between the two faces
\end{definition}
\begin{proof}
  This is the declared model definition of Has Standard Water Gravity Calibration.
\end{proof}
\begin{definition}[Satisfies Hydrostatic Pressure Law]
  \label{def:physics:phyx-mini-0392:phyxminiproblems-problemphyxmini0392-satisfieshydrostaticpressurelaw}
  \lean{PhyXMiniProblems.ProblemPhyXMini0392.SatisfiesHydrostaticPressureLaw}
  \uses{def:physics:phyx-mini-0392:phyxminiproblems-problemphyxmini0392-lengthinmeters, def:physics:phyx-mini-0392:phyxminiproblems-problemphyxmini0392-massdensityinkilogramspercubicmeter, def:physics:phyx-mini-0392:phyxminiproblems-problemphyxmini0392-accelerationinmeterspersecondsquared, def:physics:phyx-mini-0392:phyxminiproblems-problemphyxmini0392-pressureinpascals, def:physics:phyx-mini-0392:phyxminiproblems-problemphyxmini0392-damportsetup}
  Hydrostatic equilibrium measured downward from the lake surface. On the wet face, absolute pressure is ambient pressure plus ρ g y; on the dry face, the air pressure is spatially uniform and equal to ambient pressure. This is a general governing law and contains no resultant-force value or answer choice
\end{definition}
\begin{proof}
  This is the declared model definition of Satisfies Hydrostatic Pressure Law.
\end{proof}
\begin{definition}[net Pressure In Pascals At Depth Meters]
  \label{def:physics:phyx-mini-0392:phyxminiproblems-problemphyxmini0392-netpressureinpascalsatdepthmeters}
  \lean{PhyXMiniProblems.ProblemPhyXMini0392.netPressureInPascalsAtDepthMeters}
  \uses{def:physics:phyx-mini-0392:phyxminiproblems-problemphyxmini0392-pressureinpascals, def:physics:phyx-mini-0392:phyxminiproblems-problemphyxmini0392-damportsetup}
  Net water-minus-air pressure readout at a downward depth in metres
\end{definition}
\begin{proof}
  This is the declared model definition of net Pressure In Pascals At Depth Meters.
\end{proof}
\begin{definition}[Satisfies Pressure Resultant Law]
  \label{def:physics:phyx-mini-0392:phyxminiproblems-problemphyxmini0392-satisfiespressureresultantlaw}
  \lean{PhyXMiniProblems.ProblemPhyXMini0392.SatisfiesPressureResultantLaw}
  \uses{def:physics:phyx-mini-0392:phyxminiproblems-problemphyxmini0392-lengthinmeters, def:physics:phyx-mini-0392:phyxminiproblems-problemphyxmini0392-forceinnewtons, def:physics:phyx-mini-0392:phyxminiproblems-problemphyxmini0392-damportsetup, def:physics:phyx-mini-0392:phyxminiproblems-problemphyxmini0392-netpressureinpascalsatdepthmeters}
  The resultant of a pressure traction on the vertical rectangular port. A horizontal strip of height dy has area portWidth * dy, so integrating the net pressure over the full height gives the horizontal force. This generic law contains no supplied numerical values or answer choices
\end{definition}
\begin{proof}
  This is the declared model definition of Satisfies Pressure Resultant Law.
\end{proof}
\begin{definition}[Answer Choice]
  \label{def:physics:phyx-mini-0392:phyxminiproblems-problemphyxmini0392-answerchoice}
  \lean{PhyXMiniProblems.ProblemPhyXMini0392.AnswerChoice}
  Labels printed beside the four proposed force magnitudes
\end{definition}
\begin{proof}
  This is the declared model definition of Answer Choice.
\end{proof}
\begin{definition}[displayed Force In Kilonewtons]
  \label{def:physics:phyx-mini-0392:phyxminiproblems-problemphyxmini0392-displayedforceinkilonewtons}
  \lean{PhyXMiniProblems.ProblemPhyXMini0392.displayedForceInKilonewtons}
  \uses{def:physics:phyx-mini-0392:phyxminiproblems-problemphyxmini0392-answerchoice}
  Force in kilonewtons printed beside each answer label
\end{definition}
\begin{proof}
  This is the declared model definition of displayed Force In Kilonewtons.
\end{proof}
\begin{definition}[Is Unique Closest Answer]
  \label{def:physics:phyx-mini-0392:phyxminiproblems-problemphyxmini0392-isuniqueclosestanswer}
  \lean{PhyXMiniProblems.ProblemPhyXMini0392.IsUniqueClosestAnswer}
  \uses{def:physics:phyx-mini-0392:phyxminiproblems-problemphyxmini0392-answerchoice, def:physics:phyx-mini-0392:phyxminiproblems-problemphyxmini0392-displayedforceinkilonewtons}
  A choice is strictly nearer to the ideal force than every alternative
\end{definition}
\begin{proof}
  This is the declared model definition of Is Unique Closest Answer.
\end{proof}
% --- Archon declaration topology end ---
% --- Archon physics formalization source end ---
